\section{The Chern classes of a vector bundle}
\label{section-chern-classes-vector-bundles}

\noindent
We can use the projective space bundle formula to define the
Chern classes of a rank $r$ vector bundle in terms of the expansion
of $c_1(\mathcal{O}(1))^r$ in terms of the lower powers, see
formula (\ref{equation-chern-classes}).
The reason for the signs will be explained later.

\begin{definition}
\label{definition-chern-classes}
Let $(S, \delta)$ be as in Situation \ref{situation-setup}.
Let $X$ be locally of finite type over $S$.
Assume $X$ is integral and $n = \dim_\delta(X)$.
Let $\mathcal{E}$ be a finite locally free sheaf of rank $r$
on $X$. Let $(\pi : P \to X, \mathcal{O}_P(1))$ be the projective space
bundle associated to $\mathcal{E}$.
\begin{enumerate}
\item By Lemma \ref{lemma-chow-ring-projective-bundle} there are
elements $c_i \in \CH_{n - i}(X)$, $i = 0, \ldots, r$
such that $c_0 = [X]$, and
\begin{equation}
\label{equation-chern-classes}
\sum\nolimits_{i = 0}^r
(-1)^i c_1(\mathcal{O}_P(1))^i \cap \pi^*c_{r - i}
= 0.
\end{equation}
\item With notation as above we set
$c_i(\mathcal{E}) \cap [X] = c_i$
as an element of $\CH_{n - i}(X)$.
We call these the {\it Chern classes of $\mathcal{E}$ on $X$}.
\item The {\it total Chern class of $\mathcal{E}$ on $X$}
is the combination
$$
c({\mathcal E}) \cap [X] =
c_0({\mathcal E}) \cap [X]
+ c_1({\mathcal E}) \cap [X] + \ldots
+ c_r({\mathcal E}) \cap [X]
$$
which is an element of
$\CH_*(X) = \bigoplus_{k \in \mathbf{Z}} \CH_k(X)$.
\end{enumerate}
\end{definition}

\noindent
Let us check that this does not give a new notion in case the
vector bundle has rank $1$.

\begin{lemma}
\label{lemma-first-chern-class}
Let $(S, \delta)$ be as in Situation \ref{situation-setup}.
Let $X$ be locally of finite type over $S$.
Assume $X$ is integral and $n = \dim_\delta(X)$.
Let $\mathcal{L}$ be an invertible $\mathcal{O}_X$-module.
The first Chern class of $\mathcal{L}$ on $X$ of
Definition \ref{definition-chern-classes}
is equal to the Weil divisor associated to $\mathcal{L}$
by Definition \ref{definition-divisor-invertible-sheaf}.
\end{lemma}

\begin{proof}
In this proof we use $c_1(\mathcal{L}) \cap [X]$ to denote the
construction of Definition \ref{definition-divisor-invertible-sheaf}.
Since $\mathcal{L}$ has rank $1$ we have
$\mathbf{P}(\mathcal{L}) = X$ and
$\mathcal{O}_{\mathbf{P}(\mathcal{L})}(1) = \mathcal{L}$
by our normalizations. Hence (\ref{equation-chern-classes})
reads
$$
(-1)^1 c_1(\mathcal{L}) \cap c_0 + (-1)^0 c_1 = 0
$$
Since $c_0 = [X]$, we conclude $c_1 = c_1(\mathcal{L}) \cap [X]$
as desired.
\end{proof}

\begin{remark}
\label{remark-equation-signs}
We could also rewrite equation \ref{equation-chern-classes} as
\begin{equation}
\label{equation-signs}
\sum\nolimits_{i = 0}^r
c_1(\mathcal{O}_P(-1))^i \cap \pi^*c_{r - i}
= 0.
\end{equation}
but we find it easier to work with the tautological quotient
sheaf $\mathcal{O}_P(1)$ instead of
its dual.
\end{remark}




